\textbf{Phần 1. Trắc nghiệm lựa chọn}
\setcounterpageref{ex}{0}

\begin{ex} %[3D2V65]
Tính giá trị của biểu thức: $A = 35\,000 + 45\,000 - 28\,000$.
\choice
{\True $52\,000$}
{$50\,000$}
{$54\,000$}
{$48\,000$}
\loigiai{
$35\,000 + 45\,000 = 80\,000$; $80\,000 - 28\,000 = 52\,000$.
\textbf{Đáp án của câu hỏi là A}
}
\end{ex}

\begin{ex} %[3D2V65]
Tìm $x$, biết: $x + 24\,000 = 70\,000 - 16\,000$.
\choice
{$28\,000$}
{$32\,000$}
{\True $30\,000$}
{$35\,000$}
\loigiai{
$x + 24\,000 = 54\,000 \Rightarrow x = 54\,000 - 24\,000 = 30\,000$.
\textbf{Đáp án của câu hỏi là C}
}
\end{ex}

\begin{ex} %[3D2V65]
Một kho thóc có $80\,000\text{ kg}$ thóc. Ngày đầu xuất $25\,000\text{ kg}$, ngày thứ hai xuất nhiều hơn ngày đầu $5000\text{ kg}$. Trong kho còn lại bao nhiêu ki-lô-gam thóc?
\choice
{$30\,000\text{ kg}$}
{\True $25\,000\text{ kg}$}
{$20\,000\text{ kg}$}
{$35\,000\text{ kg}$}
\loigiai{
Ngày hai xuất: $25\,000 + 5000 = 30\,000\text{ kg}$. Tổng xuất: $25\,000 + 30\,000 = 55\,000\text{ kg}$. Còn lại: $80\,000 - 55\,000 = 25\,000\text{ (kg)}$.
\textbf{Đáp án của câu hỏi là B}
}
\end{ex}

\begin{ex} %[3D2V65]
Tính nhanh giá trị: $42\,500 + 17\,800 + 7500 + 2200$.
\choice
{$68\,000$}
{$72\,000$}
{$69\,000$}
{\True $70\,000$}
\loigiai{
$(42\,500 + 7500) + (17\,800 + 2200) = 50\,000 + 20\,000 = 70\,000$.
\textbf{Đáp án của câu hỏi là D}
}
\end{ex}

\begin{ex} %[3D2V65]
Tổng của hai số là 75 000, hiệu của hai số là 15 000. Số lớn là
\choice
{\True $45\,000$}
{$30\,000$}
{$40\,000$}
{$50\,000$}
\loigiai{
Số lớn: $(75\,000 + 15\,000) : 2 = 45\,000$.
\textbf{Đáp án của câu hỏi là A}
}
\end{ex}

\begin{ex} %[3D2V65]
Một mảnh đất hình chữ nhật có nửa chu vi là $450\text{ m}$, chiều dài hơn chiều rộng $50\text{ m}$. Chu vi mảnh đất đó là
\choice
{$500\text{ m}$}
{$800\text{ m}$}
{\True $900\text{ m}$}
{$1000\text{ m}$}
\loigiai{
Chu vi bằng nửa chu vi nhân 2: $450 \times 2 = 900\text{ (m)}$.
\textbf{Đáp án của câu hỏi là C}
}
\end{ex}

\begin{ex} %[3D2V65]
Một quỹ từ thiện nhận được ba đợt ủng hộ: đợt một $35\,000\,000$ đồng, đợt hai $28\,000\,000$ đồng, đợt ba $37\,000\,000$ đồng. Tổng số tiền quỹ nhận được là
\choice
{$90\,000\,000$ đồng}
{\True $100\,000\,000$ đồng}
{$95\,000\,000$ đồng}
{$110\,000\,000$ đồng}
\loigiai{
$35 + 28 + 37 = 100$ triệu đồng.
\textbf{Đáp án của câu hỏi là B}
}
\end{ex}

\begin{ex} %[3D2V65]
Tìm số tự nhiên có năm chữ số, biết rằng nếu thêm 1 đơn vị vào số đó thì được số tròn chục nghìn bé nhất có sáu chữ số ($100\,000$). Số đó là
\choice
{$99\,990$}
{$99\,000$}
{$90\,000$}
{\True $99\,999$}
\loigiai{
$100\,000 - 1 = 99\,999$.
\textbf{Đáp án của câu hỏi là D}
}
\end{ex}

\begin{ex} %[3D2V65]
Tính giá trị của biểu thức: $85\,000 - (30\,000 + 25\,000) + 12\,000$.
\choice
{\True $42\,000$}
{$40\,000$}
{$45\,000$}
{$38\,000$}
\loigiai{
$85\,000 - 55\,000 + 12\,000 = 30\,000 + 12\,000 = 42\,000$.
\textbf{Đáp án của câu hỏi là A}
}
\end{ex}

\begin{ex} %[3D2V65]
Một nhà máy có 65 000 sản phẩm. Người ta chuyển đi 3 chuyến xe, mỗi chuyến chở 15 000 sản phẩm. Trong nhà máy còn lại bao nhiêu sản phẩm?
\choice
{$15\,000$ sản phẩm}
{$25\,000$ sản phẩm}
{\True $20\,000$ sản phẩm}
{$30\,000$ sản phẩm}
\loigiai{
Đã chuyển: $15\,000 \times 3 = 45\,000$. Còn lại: $65\,000 - 45\,000 = 20\,000$ sản phẩm.
\textbf{Đáp án của câu hỏi là C}
}
\end{ex}

\begin{ex} %[3D2V65]
Số bé là 24 000, số lớn gấp 3 lần số bé. Hiệu giữa hai số là
\choice
{$72\,000$}
{\True $48\,000$}
{$50\,000$}
{$46\,000$}
\loigiai{
Số lớn: $24\,000 \times 3 = 72\,000$. Hiệu: $72\,000 - 24\,000 = 48\,000$.
\textbf{Đáp án của câu hỏi là B}
}
\end{ex}

\begin{ex} %[3D2V65]
Một cuộn dây cáp dài $100\,000\text{ m}$. Lần đầu cắt $35\,000\text{ m}$, lần sau cắt $42\,000\text{ m}$. Cuộn dây cáp còn lại dài bao nhiêu mét?
\choice
{$25\,000\text{ m}$}
{$22\,000\text{ m}$}
{$33\,000\text{ m}$}
{\True $23\,000\text{ m}$}
\loigiai{
Đã cắt: $35\,000 + 42\,000 = 77\,000\text{ m}$. Còn lại: $100\,000 - 77\,000 = 23\,000\text{ (m)}$.
\textbf{Đáp án của câu hỏi là D}
}
\end{ex}

\textbf{Phần 2. Trắc nghiệm đúng sai}
\setcounterpageref{ex}{0}

\begin{ex} %[3D2V65]
Xét tính đúng sai của các phép tính tổng hợp phạm vi 100 000:
\choiceTF
{\True Biểu thức $50\,000 + 30\,000 - 20\,000$ có giá trị bằng $60\,000$.}
{$100\,000 - 40\,000 = 50\,000$.}
{\True Tổng của hai số tròn nghìn luôn là một số tròn nghìn.}
{Hiệu giữa số lớn nhất có 5 chữ số và số bé nhất có 5 chữ số là 90 000.}
\loigiai{
\begin{itemize}
    \item Ý a) Đúng: $80\,000 - 20\,000 = 60\,000$.
    \item Ý b) Sai: Sai vì $100\,000 - 40\,000 = 60\,000$.
    \item Ý c) Đúng: Tính chất của phép cộng tròn nghìn.
    \item Ý d) Sai: Sai vì $99\,999 - 10\,000 = 89\,999$.
\end{itemize}
\begin{itemize}
    \item a) \textbf{Đúng}.
    \item b) \textbf{Sai}.
    \item c) \textbf{Đúng}.
    \item d) \textbf{Sai}.
\end{itemize}
}
\end{ex}

\begin{ex} %[3D2V65]
Một thành phố có hai quận: quận A có 45 000 dân, quận B có 55 000 dân:
\choiceTF
{\True Tổng số dân của cả hai quận là 100 000 người.}
{\True Quận B nhiều hơn quận A là 10 000 người.}
{Dân số quận B gấp đôi dân số quận A.}
{\True Cả hai số dân đều là các số tròn nghìn và tròn chục nghìn (đối với quận B là tròn chục nghìn, quận A là tròn nghìn).}
\loigiai{
\begin{itemize}
    \item Ý a) Đúng: $45\,000 + 55\,000 = 100\,000$ người.
    \item Ý b) Đúng: $55\,000 - 45\,000 = 10\,000$ người.
    \item Ý c) Sai: Sai vì $45\,000 \times 2 = 90\,000 \neq 55\,000$.
    \item Ý d) Đúng: Đặc điểm của các số.
\end{itemize}
\begin{itemize}
    \item a) \textbf{Đúng}.
    \item b) \textbf{Đúng}.
    \item c) \textbf{Sai}.
    \item d) \textbf{Đúng}.
\end{itemize}
}
\end{ex}

\textbf{Phần 3. Trả lời ngắn}
\setcounterpageref{ex}{0}

\begin{ex} %[3D2V65]
Tính giá trị biểu thức: $48\,000 + 32\,000 - 25\,000$.
\shortans{55000}
\loigiai{
\textbf{Đáp số:} $55000$\\
$80\,000 - 25\,000 = 55\,000$.
}
\end{ex}

\begin{ex} %[3D2V65]
Tìm số tự nhiên $x$, biết: $x - 18\,500 = 51\,500$.
\shortans{70000}
\loigiai{
\textbf{Đáp số:} $70000$\\
$x = 51\,500 + 18\,500 = 70\,000$.
}
\end{ex}

\begin{ex} %[3D2V65]
Tính nhanh: $27\,500 + 38\,000 + 12\,500 + 12\,000$.
\shortans{90000}
\loigiai{
\textbf{Đáp số:} $90000$\\
$(27\,500 + 12\,500) + (38\,000 + 12\,000) = 40\,000 + 50\,000 = 90\,000$.
}
\end{ex}

\begin{ex} %[3D2V65]
Một trang trại có 85 000 con gà, đã bán đi 56 000 con. Hỏi trang trại còn lại bao nhiêu con gà?
\shortans{29000}
\loigiai{
\textbf{Đáp số:} $29000$\\
$85\,000 - 56\,000 = 29\,000$ con.
}
\end{ex}

\begin{ex} %[3D2V65]
Tìm hiệu giữa 100 000 và 67 800.
\shortans{32200}
\loigiai{
\textbf{Đáp số:} $32200$\\
$100\,000 - 67\,800 = 32\,200$.
}
\end{ex}

\begin{ex} %[3D2V65]
Tổng hai số là 90 000, số bé là 35 000. Tìm số lớn.
\shortans{55000}
\loigiai{
\textbf{Đáp số:} $55000$\\
Số lớn là: $90\,000 - 35\,000 = 55\,000$.
}
\end{ex}
